\chapter[Humanoids and Whole-Body VLA]{Humanoids and\\Whole-Body VLA}

\section*{학습 목표}

이 장은 humanoid VLA를 arm-only manipulation VLA의 단순 확장으로 보지 않는다. humanoid에서는 torso, head, arms, hands, legs, balance, locomotion, gaze, bimanual coordination, whole-body contact가 하나의 closed-loop system 안에서 결합된다. 따라서 VLA는 semantic goal을 물리적으로 실행 가능한 whole-body objective로 변환해야 하고, whole-body controller는 그 objective를 balance와 contact constraint 안에서 실행해야 한다.

\begin{enumerate}
  \item humanoid VLA가 arm-only VLA와 근본적으로 다른 이유를 설명한다.
  \item whole-body state, action, constraint를 정식화한다.
  \item VLA policy와 whole-body controller의 책임 경계를 정의한다.
  \item human video, teleoperation, simulation이 humanoid data engine에서 맡는 역할을 구분한다.
\end{enumerate}

\section{Humanoid VLA is not high-dimensional arm control}

tabletop manipulator에서는 end-effector pose와 gripper command만으로 task를 상당 부분 표현할 수 있다. humanoid에서는 같은 ``pick up'' instruction도 base pose, foot placement, torso lean, arm reachability, gaze, hand posture, balance margin을 함께 결정해야 한다.

arm-only VLA를 다음과 같이 쓰자.
\begin{equation}
  \act_t^{arm}
  =
  (
  \Delta x_{ee,t},
  \Delta R_{ee,t},
  g_t
  ).
  \label{eq:ch20-arm-action}
\end{equation}
humanoid whole-body action은 훨씬 넓다.
\begin{equation}
  \act_t^{hum}
  =
  (
  \act_t^{base},
  \act_t^{torso},
  \act_t^{head},
  \act_t^{left\ arm},
  \act_t^{right\ arm},
  \act_t^{hands},
  \act_t^{legs}
  ).
  \label{eq:ch20-humanoid-action}
\end{equation}
그러나 핵심은 dimension 수가 아니라 coupling이다. arm target이 balance를 바꾸고, foot placement가 reachability를 바꾸며, gaze가 perception quality를 바꾸고, hand contact가 locomotion stability를 바꾼다.

\begin{definitionbox}{Whole-body VLA}
Whole-body VLA는 language-conditioned policy가 whole-body objective, contact mode, locomotion-manipulation coordination, gaze, hand task를 함께 결정하고, downstream whole-body controller가 balance와 physical constraint 안에서 이를 실행하는 robot policy stack이다.
\end{definitionbox}

\section{Whole-body state}

humanoid state는 joint configuration만으로 충분하지 않다. base pose, contact state, center of mass, support region, object relation, human proximity가 포함되어야 한다.
\begin{equation}
  \state_t^{hum}
  =
  (
  \conf_t,
  \dot{\conf}_t,
  x_{base,t},
  x_{com,t},
  \mathcal{C}_t,
  \mathcal{O}_t,
  h_t
  ).
  \label{eq:ch20-humanoid-state}
\end{equation}
$\mathcal{C}_t$는 contact set, $\mathcal{O}_t$는 object state, $h_t$는 human/context state이다.

belief state도 중요하다.
\begin{equation}
  \bel_t^{hum}
  =
  P(
  \state_t^{hum}
  \mid
  \obs_{\leq t},
  \ctrlcmd_{<t}
  ).
  \label{eq:ch20-humanoid-belief}
\end{equation}
humanoid는 occlusion이 많고 self-motion이 크다. head, arm, torso가 움직이면 camera extrinsic과 visible region이 계속 바뀐다. 따라서 gaze control과 state estimation은 action policy의 일부가 된다.

\section{Balance and support constraints}

humanoid control에서 balance는 선택 기능이 아니다. center of mass projection이 support polygon 안에 있거나, dynamic balance condition을 만족해야 한다.
\begin{equation}
  \Pi_{xy}(x_{com,t})
  \in
  \mathcal{P}_{support}
  \quad
  \text{for quasi-static tasks}.
  \label{eq:ch20-support-polygon}
\end{equation}
dynamic task에서는 momentum과 contact wrench까지 봐야 한다.
\begin{equation}
  \dot{h}_{mom}
  =
  \sum_{i \in \mathcal{C}_t}
  J_i(\conf_t)^{\top} \lambda_i
  +
  \tau_{ext}.
  \label{eq:ch20-momentum}
\end{equation}
여기서 $\lambda_i$는 contact wrench, $J_i$는 contact Jacobian이다.

VLA가 whole-body joint command를 직접 낸다면 이런 constraint를 모두 내부화해야 한다. 실제로는 VLA가 objective를 내고 whole-body controller가 constraint를 만족시키는 구조가 더 안전하다.

\section{Interface to whole-body controller}

VLA와 whole-body controller 사이의 interface는 다음 tuple로 정의할 수 있다.
\begin{equation}
  z_t^{wbc}
  =
  (
  x_{ee}^{L},
  x_{ee}^{R},
  x_{gaze},
  x_{base},
  \mathcal{C}_{des},
  \mathcal{W}_{task},
  \mathcal{C}_{safe}
  ).
  \label{eq:ch20-wbc-interface}
\end{equation}
왼손과 오른손 목표, gaze target, base target, desired contact mode, task weight, safety constraint가 포함된다.

whole-body controller는 보통 다음 quadratic program 형태로 쓸 수 있다.
\begin{equation}
  \begin{aligned}
  \min_{\ddot{\conf},\tau,\lambda}
  \quad&
  \sum_{k}
  \|J_k(\conf)\ddot{\conf}+\dot{J}_k\dot{\conf}-\ddot{x}_k^{des}\|_{W_k}^{2}
  +
  \|\tau\|_{W_{\tau}}^{2}
  \\
  \mathrm{s.t.}\quad&
  M(\conf)\ddot{\conf}+b(\conf,\dot{\conf})
  =
  S^{\top}\tau + J_c^{\top}\lambda,
  \\
  &
  \lambda \in \mathcal{K}_{friction},
  \qquad
  \conf \in \mathcal{Q}_{limit},
  \qquad
  \tau \in \mathcal{T}_{limit}.
  \end{aligned}
  \label{eq:ch20-wbc-qp}
\end{equation}
VLA의 역할은 이 QP를 대체하는 것이 아니라, 어떤 task objective와 contact mode를 줄지 결정하는 것이다.

\begin{table}[h]
  \centering
  \caption{VLA policy와 whole-body controller의 책임 경계}
  \begin{tabularx}{\linewidth}{p{0.24\linewidth}X X}
    \toprule
    구성 요소 & VLA가 결정하기 좋은 것 & Controller가 보장해야 할 것 \\
    \midrule
    Base & 어디로 이동할지, 어느 방향을 볼지 & footstep feasibility, collision, stability \\
    Arms & 어느 손으로 어떤 object를 다룰지 & IK, joint limit, smooth tracking \\
    Hands & grasp type, release timing & force limit, contact stability \\
    Gaze & 무엇을 봐야 하는지 & camera motion smoothness, calibration \\
    Contact & 기대 contact phase & friction cone, wrench feasibility \\
    Safety & task-level 금지 조건 & torque, speed, distance hard limit \\
    \bottomrule
  \end{tabularx}
  \label{tab:ch20-vla-wbc}
\end{table}

\section{Locomotion-manipulation coupling}

humanoid는 조작하기 위해 걸어야 하고, 걷는 동안 조작 대상과 시야를 유지해야 한다. 따라서 locomotion과 manipulation은 분리된 문제가 아니다.
\begin{equation}
  \begin{aligned}
  \max_{\xi_{base}, \xi_{arm}}\quad
  &R_{task}(\xi_{base},\xi_{arm},\ell)\\
  &-\lambda_{bal} C_{balance}(\xi_{base},\xi_{arm})\\
  &-\lambda_{col} C_{collision}(\xi_{base},\xi_{arm}).
  \end{aligned}
  \label{eq:ch20-loco-manip}
\end{equation}
base trajectory $\xi_{base}$와 arm trajectory $\xi_{arm}$는 함께 최적화되어야 한다. 예를 들어 선반 위 물체를 잡는 task에서 base가 너무 가까우면 arm collision이 생기고, 너무 멀면 reachability가 부족하다.

\begin{warningbox}{Arm-first planning의 한계}
humanoid에서 arm trajectory를 먼저 만들고 나중에 base와 balance를 맞추려 하면 실패가 많아진다. reachability, visibility, support polygon은 task planning 초기에 들어가야 한다.
\end{warningbox}

\section{Bimanual and dexterous control}

bimanual manipulation은 두 팔을 독립적으로 제어하는 것이 아니다. 양손의 역할 분담, force closure, object frame, relative motion이 필요하다.
\begin{equation}
  x_{rel}
  =
  (x_{ee}^{L})^{-1} x_{ee}^{R},
  \qquad
  x_{obj}
  =
  \mathcal{G}(x_{ee}^{L},x_{ee}^{R},c_L,c_R).
  \label{eq:ch20-bimanual}
\end{equation}
한 손은 fixture 역할을 하고 다른 손은 manipulation을 할 수 있다. 또는 두 손이 같은 object를 들고 force를 분배해야 할 수 있다.

dexterous hand는 더 어렵다. action이 finger joint command로 내려가면 contact state가 급격히 늘어난다.
\begin{equation}
  c_t^{hand}
  =
  \{
  (i,j,n_{ij},f_{ij})
  \mid
  \mathrm{finger}\ i
  \;\mathrm{contacts}\;
  \mathrm{object}\ j
  \}.
  \label{eq:ch20-hand-contact}
\end{equation}
VLA가 직접 finger-level command를 내는 것보다 grasp primitive, in-hand manipulation objective, tactile feedback policy를 결합하는 방식이 현실적일 수 있다.

\section{Dual-system architecture}

humanoid VLA는 빠른 visuomotor action과 느린 embodied reasoning을 분리하는 dual-system 구조가 자연스럽다.
\begin{equation}
  z_t^{high}
  =
  \rho_{\psi}
  (\ell,\obs_{\leq t},\bel_t^{hum},m_t),
  \qquad
  z_t^{wbc}
  =
  \pi_{\theta}
  (\obs_{\leq t},z_t^{high}).
  \label{eq:ch20-dual-system}
\end{equation}
상위 reasoning은 task, subgoal, tool use, recovery를 담당한다. 하위 VLA는 whole-body objective를 생성한다. WBC와 locomotion controller는 physics constraint를 만족시킨다.

GR00T N1은 이 관점을 읽기 좋은 최신 사례다 \cite{grootn12025}. 해당 계열은 vision-language module이 환경과 instruction을 해석하고, diffusion transformer motor module이 real-time motor action을 생성하는 dual-system humanoid VLA 구조를 제안한다. 또한 real robot trajectories, human videos, synthetic datasets를 섞는 data mixture를 사용한다. Real VLA 관점에서 중요한 점은 humanoid foundation model을 단일 policy head로 보지 않고, semantic interpretation, motor generation, data mixture, deployment embodiment를 함께 봐야 한다는 것이다.

\begin{warningbox}{GR00T N1을 읽을 때 확인할 것}
Humanoid VLA claim은 ``generalist humanoid''라는 이름만으로 평가하면 안 된다. vision-language module과 motor module의 handoff, action frequency, whole-body controller interface, balance constraint, human-video와 robot-trajectory data의 mixture ratio, real robot evaluation task를 함께 확인해야 한다.
\end{warningbox}

\begin{casebox}{Case study: GR00T N1을 humanoid system contract로 읽기}
GR00T N1의 중요한 점은 humanoid foundation model을 단일 거대 policy로만 제시하지 않는다는 데 있다 \cite{grootn12025}. vision-language module과 diffusion transformer motor module의 분리는 Chapter 18의 hierarchy, Chapter 19의 real-time budget, Chapter 20의 safety envelope와 직접 연결된다.

\begin{center}
\small
\begin{tabularx}{\linewidth}{@{}p{0.2\linewidth}X X@{}}
\toprule
계약 항목 & 읽어야 할 내용 & Real VLA 질문 \\
\midrule
High-level module & instruction, scene, memory, subgoal 처리 & semantic plan이 언제 motor objective로 바뀌는가? \\
Motor module & diffusion transformer 기반 action generation & action frequency와 WBC interface가 명시되는가? \\
Data mixture & robot trajectory, human video, synthetic data & embodiment mismatch와 retargeting error를 어떻게 다루는가? \\
Safety/eval & whole-body task와 human proximity & balance failure, near miss, intervention이 보고되는가? \\
\bottomrule
\end{tabularx}
\end{center}

이 사례의 결론은 ``humanoid도 VLA가 된다''가 아니다. 더 정확한 결론은 humanoid VLA일수록 model claim이 controller, balance, data mixture, safety case와 더 강하게 묶인다는 것이다.
\end{casebox}

\begin{table}[h]
  \centering
  \caption{Humanoid VLA stack}
  \begin{tabularx}{\linewidth}{p{0.22\linewidth}X X}
    \toprule
    Layer & 입력 & 출력 \\
    \midrule
    Reasoning & language, scene, memory & subgoal, hand choice, recovery mode \\
    VLA policy & images, proprioception, subgoal & WBC target, gaze, contact mode \\
    Whole-body controller & target, state, constraint & torque, velocity, acceleration command \\
    Locomotion controller & base target, terrain, balance & footstep, gait, base motion \\
    Safety monitor & state, human proximity, force & slow, stop, fallback, estop \\
    \bottomrule
  \end{tabularx}
  \label{tab:ch20-stack}
\end{table}

이 구조에서 가장 중요한 것은 handoff이다. reasoning이 ``상자를 내려라''라고 말하는 것과 VLA가 ``오른손으로 잡고 왼발을 앞으로 내딛어라''라고 구체화하는 것 사이에는 많은 state abstraction이 필요하다.

\section{Humanoid data engine}

humanoid data는 arm-only manipulation보다 훨씬 비싸다. human video, teleoperation, simulation, robot autonomous rollout이 서로 다른 역할을 맡는다.

\begin{table}[h]
  \centering
  \caption{Humanoid VLA data source}
  \begin{tabularx}{\linewidth}{p{0.22\linewidth}X X}
    \toprule
    Source & 제공하는 것 & 한계 \\
    \midrule
    Human video & semantic prior, body-object relation, long-horizon script & embodiment mismatch, missing force/proprioception \\
    Teleoperation & robot-valid action, contact, recovery example & cost, operator bias, limited scale \\
    Simulation & rare failure, safety stress, large variation & sim-to-real gap, contact fidelity \\
    Autonomous rollout & actual policy distribution & unsafe exploration, reset cost \\
    Motion capture & whole-body kinematic prior & robot dynamics와 actuator limit mismatch \\
    \bottomrule
  \end{tabularx}
  \label{tab:ch20-data}
\end{table}

human video를 robot action으로 쓰려면 retargeting이 필요하다.
\begin{equation}
  \xi_{robot}^{*}
  =
  \arg\min_{\xi}
  d(
  \mathcal{R}_{robot}(\xi),
  \mathcal{R}_{human}(\xi_{human})
  )
  +
  \lambda C_{feas}(\xi).
  \label{eq:ch20-retargeting}
\end{equation}
retargeting은 imitation만이 아니라 feasibility projection이다. human은 넘어지지 않고 할 수 있는 motion이라도 robot morphology와 actuator limit에서는 불가능할 수 있다.

\section{Safety envelope}

humanoid는 사람 크기의 mass와 reach를 갖기 때문에 safety envelope가 더 엄격해야 한다.
\begin{equation}
  \begin{aligned}
  \mathcal{E}_{safe}^{hum}
  =
  \{&
  \state,\ctrlcmd
  \mid
  d_{human} \geq d_{min},
  \|v_{ee}\| \leq v_{max},\\
  &
  \|F_{contact}\| \leq F_{max},
  m_{balance} \geq m_{min}
  \}.
  \end{aligned}
  \label{eq:ch20-safety-envelope}
\end{equation}
$m_{balance}$는 balance margin이다. manipulation success가 높아도 balance margin이 낮거나 human proximity constraint를 자주 위반하면 humanoid VLA는 배치할 수 없다.

\section{Evaluation}

humanoid VLA 평가는 task success만으로 부족하다.
\begin{equation}
  \begin{aligned}
  M_{hum}
  =
  (&
  R_{task},
  R_{balance},
  R_{fall},
  R_{contact},\\
  &
  R_{bimanual},
  R_{handover},
  T_{lat},
  R_{intervention}
  ).
  \end{aligned}
  \label{eq:ch20-humanoid-metrics}
\end{equation}
여기서 $R_{balance}$는 balance constraint 만족률, $R_{fall}$은 fall 또는 near-fall rate, $R_{contact}$는 unsafe contact rate, $R_{handover}$는 human handover safety를 뜻한다.

평가 protocol에는 다음 항목이 필요하다.
\begin{enumerate}
  \item robot morphology, actuator limit, controller rate
  \item task environment와 human proximity condition
  \item fall, near-fall, unsafe contact definition
  \item teleoperation 또는 human assistance 여부
  \item object mass, size, temperature, fragility
  \item balance margin, foot slip, self-collision log
  \item recovery behavior와 emergency stop event
\end{enumerate}

\section{Design checklist}

humanoid VLA를 설계할 때는 다음 순서가 실용적이다.
\begin{enumerate}
  \item VLA가 joint command를 직접 낼지, WBC objective를 낼지 정한다.
  \item manipulation task가 base motion과 footstep planning을 요구하는지 먼저 본다.
  \item gaze와 perception action을 action space에 포함할지 정한다.
  \item bimanual task에서는 object frame과 relative hand motion을 명시한다.
  \item teleoperation data에는 balance, contact, recovery label을 함께 저장한다.
  \item human video prior는 robot feasibility projection을 거친다.
  \item safety monitor는 human proximity, balance margin, contact force를 high rate로 감시한다.
\end{enumerate}

\chapterwork
{GR00T N1과 $\pi_{0.5}$를 humanoid dual-system, data mixture, whole-body controller interface 관점에서 읽어라 \cite{grootn12025,pi052025}.}
{humanoid VLA에서 manipulation action과 locomotion/balance action을 분리해서 평가하기 어려운 이유는 무엇인가?}
{humanoid handover task에 대해 high-level module, motor module, WBC, balance constraint, human-proximity safety envelope을 포함한 system contract를 설계하라.}

\section{요약}

Humanoid VLA는 VLA의 모든 어려움이 한 embodiment에 모인 형태다. state estimation, contact, action representation, hierarchical reasoning, real-time inference, safety, data scaling이 동시에 필요하다. 핵심은 VLA가 모든 joint command를 직접 예측하는 것이 아니라, semantic goal을 whole-body objective로 바꾸고, controller가 physics constraint 안에서 실행하도록 하는 책임 경계를 세우는 것이다.

다음 장에서는 책 전체를 마무리하면서 Real VLA가 앞으로 어떤 방향으로 발전해야 하는지 정리한다. 중요한 질문은 더 큰 모델 하나가 모든 문제를 해결하느냐가 아니라, data, embodiment, action interface, evaluation, safety가 함께 어떻게 성숙하느냐다.
